\chapter{Conclusion}

\section{Summary of Key Findings}

This study estimated the effect of ITN use on malaria infection among children under five in Kenya using data from the 2015 and 2020 Kenya Malaria Indicator Surveys. The analytical sample comprised 6,544 children across 528 clusters. We employed six causal inference methods: traditional logistic regression, Generalized Linear Mixed Models (GLMM), propensity score matching (PSM), inverse probability of treatment weighting (IPTW), doubly robust (DR) estimation, and Double Machine Learning (DML).

The key findings are:

\begin{enumerate}
	\item \textbf{Confounding is substantial and negative}: Unadjusted models showed no significant association (all p > 0.05). After adjusting for demographics and household factors, ITN users appeared to have significantly higher risk (OR > 1). Only after full environmental adjustment did a protective effect emerge in the combined sample (OR = 0.777, 95\% CI: 0.605-0.999, p = 0.049). This demonstrates that environmental factors were masking the true protective effect.
	
	\item \textbf{Clustering matters substantially}: The intra-class correlation of 50-58\% indicates that more than half of the variation in malaria risk is between clusters. Following \citet{killip2004intracluster}, ICC values above 0.05 justify the use of random effects models. The caterpillar plot confirmed substantial variation in cluster-specific malaria risk.
	
	\item \textbf{GLMM shows strong cluster-specific protective effects}: GLMM estimates showed OR = 0.557 (2015), 0.562 (2020), and 0.588 (combined), representing 41-44\% reduction in odds within a typical cluster. All were statistically significant.
	
	\item \textbf{GLMM-based causal methods reveal significant conditional effects}: PSM with caliper = 0.05 produced OR = 0.392 (2015, 61\% reduction) and 0.466 (2020, 53\% reduction). IPTW with weight truncation at 4 produced OR = 0.482-0.511 (48-52\% reduction). All were statistically significant.
	
	\item \textbf{Marginal effects are smaller}: DR produced near-null ORs (0.994, 0.992), representing the marginal population-average effect. DML using XGBoost produced OR = 0.972 (95\% CI: 0.960-0.984, p < 0.001) for the combined dataset, representing a 2.8\% reduction in malaria odds at the population level.
	
	\item \textbf{E-value indicates moderate robustness}: The E-value of 3.05 for the primary XGBoost estimate means that an unmeasured confounder would need to be associated with both ITN use and malaria by a risk ratio of at least 3.05 to fully explain away the observed effect. Given that measured confounders have RRs of 1.5-3.0, the effect is reasonably robust.
	
	\item \textbf{Heterogeneity exists}: ITN effectiveness varied significantly by wealth, residence, age, and rainfall. Strongest effects were observed in the poorer wealth quintile (OR = 0.919, p = 0.001), urban areas (OR = 0.974, p = 0.003), older children aged 36-59 months (OR = 0.963, p < 0.001), and medium to high rainfall regions (OR = 0.950-0.967).
	
	\item \textbf{Robustness confirmed}: Five alternative machine learning learners all produced ORs between 0.968 and 0.980, all statistically significant, confirming that DML estimates are robust to different specifications.
\end{enumerate}

\section{Answer to the Research Question}

The primary research question asked: What is the effect of insecticide-treated net (ITN) use on malaria infection among children under five in Kenya?

Based on the triangulation of evidence across six causal methods, we conclude that ITN use has a protective effect on malaria infection among children under five in Kenya. The magnitude of the effect depends on the estimand:

\begin{itemize}
	\item \textbf{Cluster-specific (conditional) effects} (GLMM, PSM, IPTW): 44-61\% reduction in odds. This answers: "For a child in a typical Kenyan cluster, how much does ITN use reduce their malaria risk?"
	
	\item \textbf{Population-average (marginal) effects} (DML): 2.8\% reduction in odds (approximately 2.2 percentage point reduction in absolute risk). This answers: "If we gave every child in Kenya an ITN, what would be the average reduction in malaria risk across the country?"
\end{itemize}

The protective effect is strongest for older children (36-59 months), in urban areas, in high-rainfall regions, and for the "Poorer" wealth quintile. At the population level, a 2.8\% reduction in malaria odds translates to approximately 125,000 malaria cases averted annually among Kenyan children under five (population approximately 5 million), representing a substantial public health impact.

\section{Limitations Acknowledged}

Several limitations must be acknowledged:

\begin{enumerate}
	\item \textbf{Unmeasured confounding}: Conditional ignorability is unverifiable. We included a comprehensive set of covariates, but unmeasured factors (genetic resistance, net condition, health-seeking behavior) may still confound the relationship. The E-value of 3.05 suggests moderate robustness.
	
	\item \textbf{Community effects not identifiable}: With cross-sectional data, personal and community protection cannot be separated. The DML estimate of 2.8\% likely mixes both mechanisms. From \citet{hawley2003community}, we know that community effects require coverage > 50\%. Without coverage data, we cannot estimate this threshold.
	
	\item \textbf{Positivity}: Some treated units have PS > 0.95 (3.3\%), some control units have PS < 0.05 (11.3\%), suggesting near-positivity violations. However, all treated units could be matched to control units in the overlap region.
	
	\item \textbf{SUTVA violation}: Herd immunity likely violates no-interference assumptions. One child's ITN use affects mosquito populations and thus affects other children's risk. Our estimates may be conservative.
	
	\item \textbf{Cross-sectional design}: Temporal ordering cannot be established for all covariates. We assume that environmental and socioeconomic characteristics precede ITN use and malaria infection.
	
	\item \textbf{Generalizability}: Findings are specific to Kenya and may not generalize to other transmission settings.
\end{enumerate}

\section{Main Contributions}

\subsection{Methodological Contributions}

\begin{enumerate}
	\item \textbf{Distinction between conditional and marginal effects}: We explicitly estimated both types of effects, demonstrating that they differ substantially when ICC is high (50-58\%). This clarifies the interpretation of previous studies that reported only one type of effect.
	
	\item \textbf{GLMM for propensity scores}: Most studies estimate propensity scores using logistic regression, ignoring clustering. We extended this to GLMM, accounting for the hierarchical structure of the KMIS data. This is an important methodological advancement for clustered survey data.
	
	\item \textbf{Cluster-aware cross-fitting for DML}: Standard cross-fitting can split children from the same cluster across different folds, creating information leakage. We implemented cluster-aware cross-fitting where entire clusters are assigned to the same fold.
	
	\item \textbf{Triangulation of six methods}: The convergence of all methods toward protective effects provides robust evidence that ITNs reduce malaria infection.
	
	\item \textbf{Integration of high-resolution environmental data}: We linked survey clusters to DHS Geocov environmental data using GPS coordinates, guided by a DAG-based approach to confounder selection.
\end{enumerate}

\subsection{Policy Contributions}

\begin{enumerate}
	\item \textbf{Rigorous evidence that ITNs reduce malaria}: This study provides some of the most rigorous observational evidence to date on ITN effectiveness in Kenya, using state-of-the-art causal inference methods.
	
	\item \textbf{Subgroup analysis for targeting}: The heterogeneity analysis identifies subgroups where ITNs are most effective: poorer quintile (not poorest), urban areas, older children, and high-rainfall regions. This informs targeted distribution strategies.
	
	\item \textbf{Documentation of declining ITN use}: ITN use declined from 58.0\% in 2015 to 50.5\% in 2020. This decline occurred despite continued distribution campaigns, suggesting that distribution alone is insufficient and behavior change communication is needed.
	
	\item \textbf{E-value for policy}: The E-value of 3.05 provides policymakers with a quantitative assessment of robustness to unmeasured confounding.
\end{enumerate}

\section{Future Research Directions}

Building on this work, future research should address the following questions:

\begin{enumerate}
	\item \textbf{Longitudinal studies}: How does the protective effect of ITNs change over time as nets age and lose insecticidal activity? Longitudinal data with repeated measures of net condition and malaria status are needed.
	
	\item \textbf{Community effects and coverage thresholds}: What is the minimum coverage needed to generate detectable community effects in different transmission settings? Cluster-randomized trials with varying coverage levels are needed.
	
	\item \textbf{Combination interventions}: What is the combined effect of ITNs with other interventions such as indoor residual spraying, larval source management, or vaccines? Factorial designs could estimate interactions.
	
	\item \textbf{Mechanisms of the wealth gradient}: Why are ITNs less effective for the poorest children? Is it housing quality, net age, overcrowding, or other factors? Data on housing characteristics, net age, and household crowding would help.
	
	\item \textbf{Generalizability}: Do these findings generalize to other malaria-endemic countries in Africa with different transmission intensities?
	
	\item \textbf{Resistance monitoring}: How does insecticide resistance affect ITN effectiveness over time? Entomological surveillance of mosquito resistance is needed.
	
	\item \textbf{DML variance estimation}: Following \citet{wu2025double}, future research should employ joint inference or bootstrap procedures to address the lack of double robustness in influence-function-based variance estimators for DML.
	
	\item \textbf{Positivity assessment}: \citet{westreich2010invited} recommend developing more sophisticated methods for detecting and handling positivity violations in high-dimensional covariate spaces.
\end{enumerate}

\section{Final Remarks}

This thesis demonstrates that ITN use reduces malaria infection among children under five in Kenya. The triangulation of evidence across six causal methods provides robust support for this conclusion. The distinction between cluster-specific effects (44-61\% reduction) and population-average effects (2.8\% reduction) highlights the importance of understanding which estimand is relevant for different policy questions.

From a policy perspective, the findings support continued investment in ITN distribution programs, but with attention to targeting and complementary interventions for the most vulnerable populations. The decline in ITN use from 2015 to 2020 is concerning and suggests that distribution alone is insufficient; behavior change communication is needed to promote consistent ITN use.

Methodologically, this study illustrates the value of modern causal machine learning methods for estimating treatment effects from observational data. The combination of traditional methods (logistic, GLMM) and modern methods (PSM, IPTW, DR, DML) allowed us to triangulate evidence and distinguish between conditional and marginal effects.

Finally, a key insight from this study is that the modest 2.8\% population - average effect estimated by DML does not imply that ITNs are ineffective. Rather, it reflects the fact that in many parts of Kenya, coverage is too low to generate community effects, and in high-transmission areas, baseline risk is so high that even perfect ITN use cannot completely block exposure. If coverage were universal (>80$\%$) across Kenya, the population-average effect would likely be substantially larger - possibly approaching the conditional effect of 44-61$\%$.

In summary, ITNs work. They reduce malaria infection in children. But they work better for some children than others. Targeting ITN distribution to high-risk areas and populations, combined with complementary interventions for the most vulnerable, can maximize the impact of malaria control programs in Kenya and beyond. The convergence of evidence from traditional and machine learning methods, all pointing toward a protective effect, provides confidence that continued investment in ITN distribution remains a sound public health strategy.